\section{Unified Problem Formulation}
\label{sec:problem}

\subsection{Evidence Allocation under Partial Observability}

Let $i$ index a PE sample with binary label $y_i\in\{0,1\}$, where one denotes
generic malware. Static representation $x_i^{\mathrm{s}}$ is available before
execution and produces score $p_i^{\mathrm{s}}$. A monitored run produces an
ordered dynamic sequence
$z_i=(z_{i1},\ldots,z_{iT_i})$. At prefix $k$, dynamic model $g$ produces
$p_{ik}^{\mathrm{d}}=g(z_{i1:k})$. The online task is to map the available
evidence to a fused score
\begin{equation}
p_{ik}^{\mathrm{f}}=
F\!\left(p_i^{\mathrm{s}},p_{i1:k}^{\mathrm{d}};\theta_F\right),
\label{eq:online-object}
\end{equation}
and a decision at a validation-selected threshold. Because $k$ counts API
events, it is an observation index rather than elapsed seconds.

The development task has a different resource constraint. Given labeled set
$\mathcal{L}_r$, unlabeled pool $\mathcal{U}_r$, and batch budget $b$, model
$f_r$ ranks pool samples with acquisition function $Q_r$. The oracle reveals
labels only for selected batch $\mathcal{B}_r$:
\begin{equation}
\mathcal{B}_r=\operatorname*{arg\,top}_{\substack{
\mathcal{B}\subset\mathcal{U}_r\\|\mathcal{B}|=b}}
\sum_{x\in\mathcal{B}}Q_r(x),
\qquad
\mathcal{L}_{r+1}=\mathcal{L}_r\cup\mathcal{B}_r.
\label{eq:batch-object}
\end{equation}
The retained experiment uses a pointwise top-$b$ rule; it does not optimize
joint diversity.

The audit task operates on a fixed fitted model and explanation
$\boldsymbol{\phi}_i$. It tests whether coordinates ranked highly by
$|\boldsymbol{\phi}_i|$ produce greater model-performance degradation under a
specified replacement intervention than equal-count random coordinates. It also
measures ranking agreement across fitted models. This task estimates relative
model reliance and sensitivity, not causal behavior of the executable.

These three tasks share an allocation principle but not one objective function.
Online fusion allocates weight between evidence channels; active learning
allocates labeling budget; explanation audit allocates confidence to a ranked
attribution. Reporting one metric as proof of all three would conflate distinct
estimands. We therefore use one architecture and three RQs, with separate data
and evaluation units.

\subsection{Fusion Estimands}

Fusion is assessed along three axes. ROC--AUC and average precision (AP) measure
ranking. The Brier score
\begin{equation}
\mathrm{Brier}=\frac{1}{n}\sum_{i=1}^{n}(p_i-y_i)^2
\label{eq:brier}
\end{equation}
measures squared probabilistic error. FPR and TPR measure decisions at thresholds
selected using validation data. An FPR comparison is interpreted together with
TPR because two independently chosen thresholds need not correspond to the same
operating point.

Paired differences use identical test rows. For Adaptive P versus Fixed, the
primary operating-point effects are
$\Delta\mathrm{FPR}$, $\Delta\mathrm{TPR}$, and
$\Delta\mathrm{Brier}$. For PID versus learned stacking, the report additionally
uses $\Delta\mathrm{AUC}$. Bootstrap intervals quantify paired uncertainty over
the retained cohort. The post-confirmatory fusion development used this same
cohort; its intervals describe resampling uncertainty, not external
confirmation.

\subsection{Learning-Curve Estimands}

Every active-learning run produces ROC--AUC and AP at budgets
$B\in\{200,300,\ldots,1200\}$. For metric $m(B)$, normalized area under the
learning curve is
\begin{equation}
\mathrm{AULC}(m)=\frac{1}{B_{\max}-B_{\min}}
\int_{B_{\min}}^{B_{\max}}m(B)\,dB,
\label{eq:aulc}
\end{equation}
computed by trapezoidal integration over measured budgets. This quantity rewards
useful models earlier in the acquisition process. Endpoint comparison is less
informative because every policy eventually receives the same full pool.

A full-pool reference model defines a fold-specific validation target
\begin{equation}
t_{\mathrm{fold}}=0.5+0.95\left(
\mathrm{AUC}^{\mathrm{full}}_{\mathrm{val}}-0.5\right).
\label{eq:target}
\end{equation}
Labels-to-target is the first training budget that reaches this value. It counts
queried training labels; the validation labels, original retained labels, and
full-pool reference are additional experimental resources. Low-FPR summaries
include normalized partial AUC up to FPR 5\% and TPR-AULC at FPR caps of 1\% and
5\%. These are cohort metrics, not estimates for operational FPR near
$10^{-3}$ or below.

\subsection{Explanation Object and Intervention}

For a fitted tree model and sample $x_i\in\mathbb{R}^d$, let
$\boldsymbol{\phi}_i=(\phi_{i1},\ldots,\phi_{id})$ denote signed TreeSHAP
contributions and $\phi_{i0}$ the bias. Native TreeSHAP is verified on the raw
score scale:
\begin{equation}
f_{\mathrm{raw}}(x_i)=\phi_{i0}+\sum_{j=1}^{d}\phi_{ij}.
\label{eq:additivity}
\end{equation}
This identity checks numerical reconstruction; it does not verify feature
semantics or causal correctness.

For deletion budget $k$, sample-specific mask $M_i^{(k)}$ replaces selected
coordinates with training-derived baseline $b_j$:
\begin{equation}
\widetilde{x}_{ij}=\begin{cases}
b_j,&j\in M_i^{(k)},\\
x_{ij},&\text{otherwise}.
\end{cases}
\label{eq:mask}
\end{equation}
The top mask contains the $k$ largest $|\phi_{ij}|$ values; the control contains
$k$ uniformly sampled coordinates. $b_j$ is the training-column mean or median.
The comparison estimates fitted-model sensitivity under two feature-space
interventions. A modified vector need not correspond to a buildable or
executable PE file.

\subsection{Security and Claim Boundaries}

The labels in both benchmarks denote generic benign/malware status. Neither
dataset establishes ransomware ground truth. The EMBER folds are random row
folds, not temporal or family-disjoint splits. MalBehavD-V1 dynamic sequences are
event-indexed, not wall-clock traces, and source differences may correlate with
class. The study excludes a white-box adversary who manipulates features, pool
membership, or the oracle against the complete pipeline.

The unified architecture is consequently a research framework. The fusion
component is evaluated on MalBehavD-V1; active learning and fidelity are
evaluated on EMBER. There is no measured interaction between adaptive fusion and
active learning, no end-to-end latency, and no common cohort supporting a joint
performance metric. These boundaries determine the wording of the abstract,
results, and conclusion.
